\section{Choosing an agent architecture for a servicing decision}
\label{exec13:agent-architecture}

A bank considering assisted causal analysis faces an architectural choice
before it faces a model leaderboard. It can give a language model a broad
notebook environment, surround a fixed statistical pipeline with conversational
assistance, or ask several specialist agents to negotiate an analysis. Each
choice changes which mistakes are easy to detect, what must be reviewed by a
person, and how much an apparently completed answer costs. The useful comparison
is therefore between complete analytical paths with declared scientific
obligations. Counting agent roles or successful tool calls gives little guidance
about whether a servicing decision has been answered.

Consider a manager asking: ``Should we offer an early restructuring review to
customers entering arrears?'' Let the unit be an eligible account at its first
qualifying arrears date, $X$ the information available immediately before the
offer, $A\in\{0,1\}$ the offer, and $Y$ the discounted net cash recovery over a
fixed following horizon. The offer is distinct from acceptance and from the
eventual restructuring terms. For a specified eligible population $P$, an
initial target is
\begin{equation}
 \tau_P=E_P\{Y(1)-Y(0)\}.
 \label{exec13:architecture-target}
\end{equation}
If $Y$ excludes the incremental cost $c$ of making the offer, an expected-value
decision compares $\tau_P$ with $c$. If costs are already inside $Y$, subtracting
them again is wrong. An agent should make this definition explicit before
requesting an estimator. Neither an offer-acceptance prediction nor the effect
among customers who accepted is automatically an answer to this question.

The information clock is equally consequential. A field can carry an event
date before the offer but be entered into the database after it. The admissible
record is the version available at the decision time, with the extraction and
revision rule specified. A distress assessment obtained during the offer call
cannot silently become a pretreatment covariate. Multiple accounts held by one
customer, branch policies and shared collection teams require an explicit
sampling and interference analysis; calling each row an independent customer
does not settle these issues. The simple point-treatment example below assumes
one eligible unit, no interference and a declared independent-unit law. The
longitudinal construction earlier in this chapter is needed when subsequent
contacts and evolving distress become part of the intervention.

\subsection{What the competing architectures actually contribute}

The primary literature supports several distinct proposals. Zhu and colleagues'
CausalAgent routes processing, structure learning and reporting through agents
and MCP tools. Its four-page demonstration uses the Sachs dataset; quantitative
effect estimation and counterfactual modules remain among its stated future
directions. Its conversational interface is evidence of an interaction design,
not a demonstrated solution to the servicing target in
\eqref{exec13:architecture-target}. Han and colleagues' separately authored
Causal Agent instead illustrates named graph objects in tool memory and a
generated question benchmark. Keeping these two systems distinct prevents a
result from one from being attributed to the other's implementation
\citep{v12_la_2602_11527,exec13_census_causal_agent}.

The actor/critic design of oci-agent addresses a closer statistical problem:
point-in-time observational effects under supplied unconfoundedness. It exposes
specifications, notebooks and diagnostics to a human principal. This can reduce
the effort of reproducing an analysis or finding a failed overlap diagnostic.
The critical question is whether a proposed revision preserves the target.
Removing unsupported parts of the population may make an estimate more stable
while answering a different banking question. That is a substantive decision
which an attractive report must not hide. The paper's benchmark results remain
reported source findings, and are not reproduced live-model results of this
proposal \citep{v12_la_2607_22443}.

The causal-learn+ position takes a stricter approach. Agents assist data
inspection, method explanation and interpretation around a protected discovery
core. Their proposed edges, orientations or constraints do not automatically
become scientific evidence. This reduces one path by which language-model
beliefs enter a graph, but leaves assumption elicitation, variable construction
and statistical method choice as important tasks. Its call for future
benchmarks is consistent with treating the architecture as a candidate for
evaluation rather than an established performance advantage
\citep{v12_la_2606_23608}.

ORCA emphasizes shared relational state: joins, variable roles, graphs and
estimation configurations. This is especially relevant when banking data are
spread over account, customer, contact and transaction tables. However, the
paper's conversion of a consensus PAG into an actionable DAG introduces
orientation, reversal or deletion decisions. An analysis conducted on that DAG
inherits these extra choices. The architecture should expose them as additions
to the assumptions, and preserve the original uncertain graph. The numerical
comparison is also narrower than its headline suggests: the inspected table
reports completion rates of $.83$ and $.75$, whereas the abstract claims a
42-percentage-point gain. Completion is defined as producing an ATE regardless
of its accuracy. These are reasons to preserve the source discrepancy and to
design a better denominator, not to invent a corrected experiment
\citep{exec13_census_orca}.

Two further architectures broaden the comparison. CausalSmith makes a formal
statement, proof dependencies and compiler results explicit; its own limitations
separate proof checking from unvalidated model judgments about statement
faithfulness and research quality. Model Discovery Agent proposes executable
mechanisms, then uses numerical fitting and experiment design to compare them.
The latter's model-disagreement criterion approximates information gain within
a changing candidate set. Its reported predictive tests do not establish
causal identification, and its illustrated experiments use open-loop policies.
Neither formal proof search nor mechanism proposal alone replaces the
observational assumptions needed for the servicing target
\citep{exec13_census_causalsmith,v12_la_2608_09696}.

\begin{longtable}{@{}P{.19\textwidth}P{.35\textwidth}P{.39\textwidth}@{}}
\caption{Architecture choices for a bounded servicing prototype. These are
proposal decisions based on inspected sources.}\label{exec13:architecture-choices}\\
\toprule
Design & Useful intermediate object & Proposed reuse boundary\\\midrule
\endfirsthead
\toprule Design & Useful intermediate object & Proposed reuse boundary\\\midrule
\endhead
Conversational causal tools & Named variables, graph objects and explicit
tool arguments & Adapt interaction ideas; independently check graph encoding,
tool scope and malformed-result handling.\\
Actor/critic notebooks & Analysis specification, code, diagnostics and
criticism & Adapt review artifacts; restrict execution to a frozen supported
analysis and require a new target declaration after trimming.\\
Protected discovery core & Distinct suggestions, decisions, parameters and
graph outputs & Prefer for an initial discovery assistant; retain graph
ambiguity and test the advice surrounding the core.\\
Shared relational state & Joins, populations, candidate graphs and estimators
& Adopt the concept of typed intermediate state; keep every consequential
join and orientation inside the scientific review.\\
Formal proof workflow & Exact statement, definitions, imports and proof
& Use for selected derivations; review semantic fidelity and satisfiability
separately from compilation.\\
Mechanism proposal loop & Executable mechanisms, fitted weights and experiment
candidates & Consider for a later controlled simulator; inspect model support,
action coverage and approximation error before causal interpretation.\\
\bottomrule
\end{longtable}

These alternatives suggest a restrained first prototype: one application-owned
controller, a bounded set of numerical and graph operations, and optional
language-model assistance for drafting and criticism. Multiple personas can be
an experimental arm. Their added value should be measured against the same
tools, information and total repair budget available to a simpler workflow.
This choice concerns the proposed product; it does not claim that the reviewed
external frameworks have been installed or integrated.

\subsection{A target can drift while every tool succeeds}

Suppose consistency, conditional exchangeability
$(Y(0),Y(1))\perp A\mid X$, positivity and the stated sampling law hold.
Write $e(x)=P(A=1\mid X=x)$ and $m_a(x)=E(Y\mid A=a,X=x)$.
The familiar augmented inverse-probability score is
\begin{align}
 \phi(O;m,e)={}&m_1(X)-m_0(X)
   +\frac{A\{Y-m_1(X)\}}{e(X)}\nonumber\\
 &-\frac{(1-A)\{Y-m_0(X)\}}{1-e(X)}.
 \label{exec13:architecture-score}
\end{align}
At the true nuisance functions, conditioning on $X$ removes both residual
terms, leaving $E\{\phi\mid X\}=E\{Y(1)-Y(0)\mid X\}$.
Consequently, averaging over $P$ gives $\tau_P$. Estimated nuisance functions,
cross-fitting, dependence and uncertainty require the additional arguments
developed elsewhere in this monograph. Calling this score does not establish
exchangeability or those rate conditions \citep{hernan2020}.

Now let an overlap diagnostic cause the workflow to keep only
$K=\{x:.1\leq e(x)\leq .9\}$. A uniform average of the score among retained
units targets
\begin{equation}
 \tau_K=\frac{E[1_K(X)\{m_1(X)-m_0(X)\}]}{P(X\in K)}.
 \label{exec13:architecture-trim}
\end{equation}
By contrast, define the overlap-weighted target explicitly as
\begin{equation}
 \tau_{\mathrm{ov}}=
 \frac{E[e(X)\{1-e(X)\}\{m_1(X)-m_0(X)\}]}
      {E[e(X)\{1-e(X)\}]}.
 \label{exec13:architecture-overlap}
\end{equation}
Equations \eqref{exec13:architecture-trim} and
\eqref{exec13:architecture-overlap} integrate against different measures.
They coincide only under additional restrictions, such as a constant
conditional effect; their labels cannot make them equal. If $K$ is estimated
from the same data, uncertainty about the selection rule is another issue.
The present calculation uses known propensities and a fixed rule so that this
additional difficulty does not obscure the target distinction.

In the inspected oci-agent notebook, the branch labeled ATO calls a propensity
trimming function and then applies a uniform AIPW summary to the retained
pseudo-outcomes. The serialization retains this label and exports retained
indices separately. This source path supports a trimmed-population
interpretation under the relevant assumptions; it does not implement the
weighting in \eqref{exec13:architecture-overlap}. The appropriate adaptation
therefore reads the implemented weights and retained population, not just the
display name \citep{v12_la_package_oci_agent}.

An exact generated example makes the financial consequence visible. Take
three strata with masses $(1/4,1/2,1/4)$, offer probabilities
$(.05,.5,.8)$ and conditional net-recovery effects $(-200,100,500)$ dollars.
Direct summation gives
\begin{equation}
 \tau_P=125,\qquad
 \tau_K=\frac{700}{3}\simeq233.33,\qquad
 \tau_{\mathrm{ov}}=\frac{48200}{283}\simeq170.32
 \quad\text{USD per account}.
 \label{exec13:architecture-target-numbers}
\end{equation}
Trimming retains three quarters of the original population. If the offer
cost is 180 dollars per eligible account and lies outside the recovery
outcome, the full-population expected net gain is $-55$ dollars. Substituting
the trimmed effect into that full-population calculation would reverse the
sign. A policy confined to the retained stratum is a different decision and
needs its own eligibility rule. The numbers are authored potential outcomes,
not observations about any bank. A separate 400-unit finite-population
calculation with exact nuisances reproduces all three weighted averages
\citep{exec13_architecture_exact}.

The same problem occurs in graph handoffs. Under one declared DAG, a baseline
distress measure can be a common cause of offer and recovery and belong in an
adjustment set. Under a different timing definition it can be a descendant of
the offer, so conditioning on it can remove part of the total effect. An
architecture should attach the variable definition, information time and graph
assumption to the adjustment result. Caching by column names alone cannot
preserve that meaning.

\subsection{What repeated critics can and cannot repair}

A critic is useful when it brings a distinct check: reconstructing weights,
reading a source passage, evaluating a graph criterion or reproducing a
numerical baseline. Merely repeating the same prompt with the same evidence
does not create independent scientific information. To quantify this
distinction, suppose each of three critics misses an error with marginal
probability $q=.1$. If their errors are independent, majority failure has
probability
\begin{equation}
 q^3+3q^2(1-q)=.028.
 \label{exec13:architecture-majority}
\end{equation}
Perfectly shared errors instead leave the failure probability at $.1$.
Consider an intermediate stipulated law: with probability $\gamma=.05$,
all three fail from a common cause; otherwise they fail independently with
probability $r=1/19$. Each still has marginal error
$\gamma+(1-\gamma)r=.1$, but majority failure becomes
\begin{equation}
 \gamma+(1-\gamma)\{r^3+3r^2(1-r)\}\simeq .05762.
 \label{exec13:architecture-common-error}
\end{equation}
Thus the same individual error rate can accompany quite different benefits
from voting. Exact enumeration over all eight independent-error patterns
checks the calculation. None of these probabilities is an empirical estimate
of a language model or a deployed critic \citep{exec13_architecture_exact}.

There is also an information limit. Consider two linear Gaussian worlds with
mutually independent standard normal noises in each world:
\begin{align}
 \mathcal W_1:&\quad X=U,\quad Y=\tfrac12X+\sqrt{3/4}\,V,\nonumber\\
 \mathcal W_2:&\quad Y=V',\quad X=\tfrac12Y+\sqrt{3/4}\,U'.
 \label{exec13:architecture-worlds}
\end{align}
Both observed pairs have zero mean, unit marginal variances and covariance
$1/2$, hence the same joint distribution. The effect of increasing $X$ by
one is $1/2$ in $\mathcal W_1$ and zero in $\mathcal W_2$. Suppose an
adaptive agent system receives only samples from that common law, identical
documents in both worlds and randomness independent of the world. The first
transcript step then has the same law in both worlds. If this holds through
step $t$, the next tool or agent applies the same conditional transition to
an identically distributed history, so it holds at $t+1$. Induction gives the
same distribution for the final answer. An additional critic operating on
that transcript cannot distinguish these worlds.

This argument concerns the stated information set. Genuine temporal knowledge,
an intervention or a justified scientific restriction can break its premise.
Retrieval is valuable when it supplies such information with an inspectable
provenance; fluent agreement alone cannot do so. The appropriate output before
the added information is available can be an identified set, a conditional
answer or a request to clarify the decision model. It need not be an arbitrary
scalar produced to make the workflow look complete.

Repeated repair also changes evaluation. In a deliberately simple experiment,
suppose a scientifically invalid candidate independently has probability $.05$
of passing a fallible screen on each try. Allowing twenty tries and retaining
the first apparent pass gives probability
$1-.95^{20}\simeq.6415$ of at least one false pass. Without independence, the
union bound is only $\min(1,20\times.05)=1$. This example describes independent
screening attempts, not deterministic replay of the same input, and does not
assign a statistical significance level to real agent review. It explains why
the proposed evaluation must retain all attempts, give comparators the same
repair budget and separate development feedback from held-out evaluation
\citep{exec13_architecture_exact}.

\subsection{Typed handoffs and the scope of recovery}

For the servicing prototype, a useful analysis state contains the decision
contract, data identity and clock, graph or graph class, assumptions with their
source and status, target measure, method applicability, numerical artifacts,
uncertainty statement and request history. These objects should refer to one
another by immutable identity. A revised target or graph invalidates dependent
results even when the data bytes are unchanged. A successful statistical call
must not make an unresolved assumption disappear.

The controller can enforce local invariants. An estimation request must name a
target and population; every variable must exist in the supplied schema; a
numerical interval must have finite, ordered endpoints unless the result type
explicitly supports an unbounded set; and a result must bind the method and
input versions that produced it. A graph-identification operation can check a
formal condition relative to a supplied graph. These checks support meaningful
handoffs, but none proves that the graph or observation law describes the bank.
Scientific assumption status should therefore remain visible alongside
mechanical validity.

This distinction leads to several result types with different continuation
actions. Malformed input calls for correction of the request. A missing backend
calls for an implementation or environment remedy. A mathematical
nonidentification result calls for different information or a different target.
An unsupported method calls for a scoped alternative. A valid set-valued answer
can proceed to a decision calculation that respects the set. Collapsing all of
these into either an empty successful object or a single generic failure would
discard the reason that determines the next useful action.

The inspected CausalAgent server illustrates a concrete boundary to repair
before reuse. Its OLC handler refers to a function not imported or defined in
that module and catches the resulting exception into a JSON error string.
Separately, its output helper extracts text between braces and returns an
empty dictionary on a decoding failure. These are static source findings, not
a claim that this proposal executed the upstream server. A wrapper must
validate the scientific result and propagate an explicit failure before
downstream analysis; successful transport of an error string is not a learned
graph \citep{v12_la_package_causalagent}.

Recovery has a related distinction. A graph checkpoint preserves enough state
to resume an operation; it does not automatically make external effects occur
exactly once. The inspected LangGraph interrupt restarts its node from the
beginning. An application should therefore give an expensive computation a
stable request identity, record its durable outcome and check that identity on
replay. For a remote provider call interrupted before a receipt is saved, the
outcome can remain unknown unless the provider supports lookup or idempotency.
Blind resubmission cannot be advertised as exact-once execution
\citep{v12_la04_langgraph}.

Pydantic AI supplies typed tool boundaries and usage controls, but the inspected
price check cannot enforce a monetary limit when price is unavailable. Some
usage becomes known only after a response. A prototype requiring a hard budget
needs a known pricing contract, conservative reservation and explicit behavior
for unknown cost. AutoGen's inspected save/load methods preserve model context;
they do not constitute a journal of arbitrary external side effects. Its
maintenance declaration and separate code/documentation licenses also enter
the adoption decision. These observations favor a plain controller baseline
before adding a framework whose incremental benefit has not yet been measured
\citep{v12_la04_pydantic_ai,v12_la04_autogen}.

The existing local prototype provides narrower positive evidence. Six authored
offline cases preserve blocked, unsupported, ambiguous and completed outcomes,
including directed-revision recovery and a weak moment set. The saved replay
reproduces those outcomes using real local tools with no live model calls.
This supports selected semantics and recovery behavior. It does not test
whether a language model elicits the right servicing graph, reduces human
effort or behaves correctly on an external platform. The library, CLI and MCP
contracts should carry this same distinction when transporting the result
\citep{exec13_architecture_local}.

\subsection{Selecting the next experiment by useful output and total cost}

The minimal comparison has a deterministic analyst template, a single assisted
controller and, if justified, a controller with a critic. All receive the same
documents, permitted tool implementations and task budget. Generated cases
should include correctly identified effects, ambiguity requiring a set or
refusal, misleading variable clocks, a target change after trimming and a tool
failure. A held-out task split separates adaptation of prompts and templates
from evaluation. Numerical truth comes from the declared generator, while
semantic scoring needs an independently specified reference and ambiguity
rubric. A correct refusal is a successful answer to a nonidentified task;
unnecessary refusal on an identified task is a separate error.

Let $N$ count all eligible assigned requests and $U_i$ indicate that request
$i$ produced an admissible, correct and decision-appropriate answer under the
rubric. A useful workflow summary is
\begin{equation}
 \widehat C_{\mathrm{usable}}=
 \frac{\sum_{i=1}^N(C_i^{\mathrm{tools}}+wT_i^{\mathrm{review}}
                         +C_i^{\mathrm{other}})}{\sum_{i=1}^N U_i}.
 \label{exec13:architecture-cost}
\end{equation}
It is undefined when there is no usable output. Costs for failed requests stay
in the numerator. The statistic should accompany the usable fraction, error
categories, refusal behavior, latency distribution and conditional numerical
accuracy. Reporting only accuracy among completed estimates can reward a system
that abandons difficult requests or hides configuration errors.

For planning, suppose manual review costs 15 minutes per request at 120 dollars
per hour and yields a usable fraction $.9$. Its ratio is $30/.9=33.33$ dollars.
An assisted workflow costing 3 dollars in tools plus 6 review minutes, with
usable fraction $.8$, gives $15/.8=18.75$ dollars. At 20 review minutes the
same tool and quality assumptions give $43/.8=53.75$ dollars. Equating the
assisted ratio to the manual ratio yields a review-time threshold of
$71/6\simeq11.83$ minutes. These stipulated scenarios omit fixed build costs;
they show which quantities to measure, not that assistance saves money
\citep{exec13_architecture_exact}.

The next justified experiment is consequently a bounded public/generated task
study measuring semantic failures and total review effort, with matched tools
and frozen denominators. A live-model performance claim requires a recorded
backend revision, actual calls and measured usage. A literature and architecture
proposal can proceed without that claim. Authentic bank experiments become
worth considering when the prototype isolates an uncertainty whose resolution
could change the build or investment decision. The unresolved questions are
specific: how often assumptions are elicited correctly, whether criticism adds
independent information, how much adjudication is needed and whether the
resulting gain survives a change in task population. None is answered by adding
agent roles to a diagram.
\label{exec13:agent-architecture-end}
